% Fragmento público generado desde propietarios íntegros.
% Residencia: 52.8.5.
% Fuente: ../incoming/radion_monografia_20260827/manuscrito/sections/09_sector_90120.tex:55-114
\subsubsection{Publicación espectral de la cola de órdenes \(j\ge2\)}

\paragraph{Definición del funcional espectral}

Para \(m\in\{90,120\}\), sea
\begin{equation}
L_m^{(2)}(q)=\sum_{j\ge2}\frac{q^{mj}}{j}
=-\log(1-q^m)-q^m.
\label{eq:L2}
\end{equation}
La lectura de una pareja es
\begin{equation}
T_h^{(2)}=\frac{180}{\pi}
\Bigl[
L_{90}^{(2)}(q_{h,-})-L_{90}^{(2)}(q_{h,+})
-L_{120}^{(2)}(q_{h,-})+L_{120}^{(2)}(q_{h,+})
\Bigr].
\label{eq:T-h}
\end{equation}
La reconstrucción de las seis parejas da
\begin{equation}
F_{\mathrm{spec}}=6T_h^{(2)}
=5.1499235153094574410\ldots\times10^{-8}\,{}^\circ.
\label{eq:Fspec}
\end{equation}

El acarreo del radión es
\[
F_{\mathrm{acarreo}}=\epsR^\circ
=5.1383016758575282072\ldots\times10^{-8}\,{}^\circ.
\]
No coinciden:
\begin{equation}
\chi_R:=F_{\mathrm{spec}}-F_{\mathrm{acarreo}}
=1.16218394519292338\ldots\times10^{-10}\,{}^\circ.
\label{eq:chiR}
\end{equation}

La igualdad
\[
F_{\mathrm{acarreo}}=F_{\mathrm{spec}}-\chi_R
\]
es un cambio de carta válido una vez construidas ambas salidas. No debe
presentarse como generación independiente de \(\epsR^\circ\) si \(\chi_R\) se
define por la propia resta. La cola está cerrada y es significativa; lo falso
es identificarla directamente con el acarreo.

\begin{falsifier}
Tras convertir la cola expresada en grados a la unidad angular interna,
\[
\frac{\pi}{180^\circ}F_{\mathrm{spec}}-\epsone
=2.0283936357433839\ldots\times10^{-12}\ne0.
\]
Por tanto, ni el log completo ni la cola \(j\ge2\), aun después de aplicar la
carta de unidades, son \(\epsone\). Una
segunda factorización espectral autónoma necesitaría generar \(\chi_R\) sin
usar la diferencia como definición; el cierre APP--TPK del radión no depende
de esa promoción.
\end{falsifier}

